\section{Kesimpulan}
\label{sec:kesimpulan}

Berdasarkan implementasi dan eksperimen pada tugas ini, kesimpulan yang dapat ditarik adalah sebagai
berikut.

\begin{enumerate}
	\item Kelima algoritma, yaitu KNN, DTL (C4.5), \textit{Logistic Regression} dengan regularisasi L1,
	L2, dan \textit{Elastic Net}, SVM dengan kernel \textit{linear}, \textit{polynomial}, dan
	\textit{Gaussian}, serta \textit{Naive Bayes}, berhasil diimplementasikan \textit{from scratch}.
	Setiap model dapat disimpan dan dimuat kembali.
	\item Pada protokol pembanding yang sama, KNN, \textit{Logistic Regression}, SVM, dan \textit{Naive
	Bayes} \textit{from scratch} menghasilkan Macro F1 yang sama dengan \textit{scikit-learn} sampai
	empat desimal, dengan selisih terbesar 0,0003. DTL C4.5 \textit{from scratch} lebih tinggi daripada
	konfigurasi \texttt{DecisionTreeClassifier} terbaik (0,8692 berbanding 0,8642) karena memakai
	\textit{gain ratio}, pemecahan kategorikal multicabang, penanganan nilai kosong dengan bobot
	fraksional, dan \textit{pruning}.
	\item Model terbaik adalah SVM dengan kernel yang dipelajari, dengan Macro F1 validasi silang 0,8750 (satu
	\textit{seed}) dan skor \textit{public} Kaggle 0,87912 (0,87973 untuk varian 255 kuantil). SVM dengan kernel Laplacian dan bobot fitur mendapat
	0,8703 dan 0,87464, dan DTL mendapat 0,8670 dan 0,87393.
	\item Tahap \textit{preprocessing} yang paling berpengaruh adalah pembentukan fitur ``saudara'', yaitu
	informasi dari pemesanan lain yang mirip pada data latih. Fitur \texttt{sib\_te} menjadi pemecahan
	akar pohon DTL, dan keunggulan SVM atas pohon pembanding seluruhnya berasal dari baris yang memiliki
	saudara. Sebaliknya, PCA tidak membantu SVM sehingga tidak dipakai pada model akhir.
	\item Dengan asumsi FP dua kali lebih mahal daripada FN, penanganan yang tepat adalah mengurangi
	prediksi batal: menaikkan \textit{threshold} atau memberi bobot lebih besar pada kelas tidak batal.
	Pada DTL, pembobotan kelas 2 : 1 menurunkan total biaya OOF dari 4.286 menjadi 4.124, sedangkan
	pada KNN \textit{threshold} $2/3$ menurunkan biaya \textit{hold-out} sebesar 8,87\%. Pada kedua
	model, Macro F1 turun sebagai gantinya, sehingga model untuk penginapan sebaiknya dipilih berdasarkan
	biaya, bukan hanya Macro F1.
\end{enumerate}

\section{Saran}
\label{sec:saran}

Beberapa hal yang dapat dikembangkan lebih lanjut adalah sebagai berikut.

\begin{enumerate}
	\item Kesalahan terbanyak terjadi pada baris tanpa saudara di data latih (Macro F1 sekitar 0,81,
	dibandingkan sekitar 0,94 pada baris dengan saudara). \textit{Feature engineering} lanjutan untuk
	baris tersebut, misalnya fitur dari riwayat pelanggan atau pola musiman, kemungkinan memberi
	kenaikan yang lebih besar daripada \textit{tuning} model.
	\item SVM dan KNN mahal secara komputasi sehingga pemilihan konfigurasi SVM dilakukan pada sampel
	3.000 sampai 6.000 baris. Teknik \textit{shrinking} seperti pada LIBSVM, aproksimasi kernel
	(misalnya metode Nyström), atau struktur indeks tetangga (misalnya KD-tree) dapat memungkinkan
	pemilihan konfigurasi pada seluruh data.
	\item TAN sudah lebih baik daripada \textit{Naive Bayes}, tetapi setiap fitur hanya boleh memiliki
	satu \textit{parent} selain kelas. Struktur \textit{Bayesian network} yang lebih umum, misalnya hasil
	pencarian struktur berbasis skor, dapat dieksplorasi.
	\item Validasi \textit{cost-sensitive} sebaiknya memakai biaya nyata dari penginapan, karena rasio
	2 : 1 pada tugas ini adalah asumsi. \textit{Threshold} juga perlu dikalibrasi ulang jika rasio
	pembatalan pada data baru berubah.
\end{enumerate}
